% This file is part of the modular source for Stratified Relational Workspace Theory.
% !TeX root = ../Stratified_Relational_Workspace_Theory.tex
\theorypart{XXIII}{Relation to Prior Work}

\topic{83}{Scope, Method, and Standards of Priority}

\subtopic{83.1 Purpose of the literature review}

The preceding parts construct a unified formal theory. Construction alone does
not establish priority. This part therefore asks a narrower question:
which ingredients are established, which are adaptations to a new domain,
which combinations appear to be unusual, and which claims remain open originality candidates?

The audit uses five labels:
\[
\begin{array}{ll}
\mathsf{STD} & = \text{established concept, theorem, or technique},\\
\mathsf{ADP} & = \text{established machinery adapted to relational workspaces},\\
\mathsf{SYN} & = \text{nontrivial synthesis of established components},\\
\mathsf{CAN} & = \text{candidate original contribution requiring priority review},\\
\mathsf{EMP} & = \text{empirical claim not yet established in an uncontrolled system}.
\end{array}
\]
These labels describe the status of a claim, not its importance. A standard
result may be indispensable. A synthesis may be scientifically useful without
containing a new theorem. A candidate contribution may later be narrowed or
withdrawn when a closer precedent is found.

\subtopic{83.2 Search scope and date}

This review covers sources available through August 29, 2026. It includes:
\begin{itemize}
\item extended and distributed cognition, transactive memory, shared
 intentionality, common knowledge, collective intelligence, and group
 agency;
\item blackboard systems, global-workspace architectures, team cognition,
 hybrid intelligence, and human--AI cooperation;
\item category-theoretic limits, local systems, holonomy, sheaves, \v{C}ech
 descent, connection and sheaf Laplacians, consensus, and spectral
 perturbation;
\item statistical sufficiency, causal states, information decomposition,
 structural causal models, system identification, frames, erasure
 recovery, and probabilistic process equivalence; and
\item agency, viability, control, empowerment, damage detection, repair, and
 self-maintaining systems.
\end{itemize}
The search combined exact-title queries, component-level queries, backward
reference tracing, and comparison with recent formal work. It is a scholarly
priority review, not a legal patent search and not a proof that no uncatalogued,
unpublished, or differently named precedent exists.

\subtopic{83.3 A stronger standard than name uniqueness}

Failure to find the phrase ``Stratified Relational Workspace Theory'' in prior
literature is weak evidence. Priority is determined by mathematical and
operational equivalence, not by terminology. Let \(C\) be a claim in the
paper and \(S\) a prior source. Write
\[
S\preccurlyeq C
\]
when the main content of \(C\) follows from \(S\) after admissible relabeling,
restriction, or direct specialization. A substantive novelty residual exists
only when the claim contains content not absorbed by the closure of known
sources:
\[
\operatorname{Res}(C)
=
C\setminus
\operatorname{Cl}
\left(
\bigcup_{S\in\mathcal L}S
\right).
\]
This is a conceptual audit device rather than a set-theoretic definition of
ideas. Its purpose is to force an explicit account of what remains after the
standard machinery is removed.

\subtopic{83.4 When a synthesis counts as a contribution}

A conjunction of known ingredients is not automatically new. The present
framework treats a synthesis as a candidate contribution only when it satisfies
all three conditions:
\begin{enumerate}
\item \textbf{Non-redundancy.} Removing one of the joined layers changes the
 class of systems certified or the predictions made.
\item \textbf{Formal integration.} The layers share declared objects,
 transfer maps, interventions, or invariants rather than appearing as an
 unconnected checklist.
\item \textbf{New discriminations.} The synthesis separates at least one
 positive case from a matched near miss that the individual components do
 not separate on their own.
\end{enumerate}
The anti-centralization, matched-null, non-collapse, provenance, replacement,
and repair-replay controls are important precisely because they make the
combination testable rather than merely terminological.

\subtopic{83.5 Priority claims are coordinatewise}

The analysis does not make one global novelty claim. It uses a priority
profile
\[
\boxed{
\mathbf p(C)
=
\left(
 p_{\mathrm{concept}},
 p_{\mathrm{object}},
 p_{\mathrm{theorem}},
 p_{\mathrm{certificate}},
 p_{\mathrm{experiment}}
\right).
}
\]
For example, the philosophical concept of distributed cognition may be
standard while a specific mathematical object is adapted, a theorem is a
specialization of known spectral theory, the certificate is a synthesis, and
the empirical claim is unresolved. This prevents originality in one
coordinate from being used to imply originality in every coordinate.

\topic{84}{Philosophical and Cognitive Ancestry}

\subtopic{84.1 Extended cognition}

The extended-mind tradition argues that external resources can participate in
cognitive processes when they are appropriately integrated with an agent's
activity \cite{ClarkChalmers1998}. Its critics emphasize that causal coupling
alone is too permissive and that cognition requires a principled mark or
boundary \cite{AdamsAizawa2001,Rupert2004}. Stratified Relational Workspace
Theory inherits the possibility of extracranial realization but changes the
question. It does not ask whether a notebook, model, or collaborator literally
belongs to one individual's mind. It asks whether a distributed realization
contains a persistent, causally operative, jointly reconstructible structure
that passes explicit anti-triviality gates.

Accordingly,
\[
\operatorname{ExtendedCognition}
\centernot\Longleftrightarrow
\operatorname{SRW}.
\]
An externally scaffolded process may be cognitively extended without being
relationally irreducible. Conversely, an institution or multi-agent system may
instantiate a certified relational workspace without being an extension of any
single person's cognition.

\subtopic{84.2 Distributed cognition}

Distributed cognition studies cognitive organization across people,
representations, artifacts, and environments rather than locating all relevant
processing inside one individual \cite{Hutchins1995,HollanEtAl2000}. This is
the closest philosophical ancestor of the framework's ontology. It already
supports the central observation that the unit of analysis may be a coupled
system and that external representations can transform the computations
available to that system.

The framework should therefore not claim the following ideas as new:
\begin{itemize}
\item cognition can be distributed across heterogeneous carriers;
\item artifacts can be active components of a cognitive process;
\item representational transformations can alter system-level performance; and
\item the relevant explanatory unit can exceed an individual organism.
\end{itemize}
The added ambition is mathematical and certifying. The theory specifies local
state spaces, transfer maps, compatible sections, causal sufficiency,
coalition irreducibility, persistence, replacement identity, formation
provenance, and matched nulls. Distributed cognition supplies the ontological
motivation. It does not by itself supply this complete formal certificate.

\subtopic{84.3 Transactive memory and collective intelligence}

Transactive-memory theory models the group-level organization of ``who knows
what'' and the processes by which members encode, retrieve, and coordinate
specialized knowledge \cite{Wegner1987,Lewis2003}. Collective-intelligence
research studies stable differences in group performance that are not reduced
to the intelligence of a single member \cite{WoolleyEtAl2010}. Both traditions
are directly relevant to the coalition layer.

The distinction is:
\[
\begin{aligned}
\text{transactive memory}
&\approx
\text{distributed indexing and retrieval of expertise},\\
\text{collective intelligence}
&\approx
\text{system-level task capacity},\\
\text{relational workspace}
&\approx
\text{an invariant state plus transfer, causal, and identity structure}.
\end{aligned}
\]
A transactive-memory system may still be centralized by one coordinator. A
high-performing group may still be explainable by the best individual member.
The certificate therefore requires matched coalition risks and nulls
rather than inferring relationality from group success.

\subtopic{84.4 Shared intention, common knowledge, and group agency}

Common-knowledge theory formalizes nested epistemic conditions required by
many distributed protocols \cite{HalpernMoses1990,FaginEtAl1995}. Shared-action
accounts analyze jointly intended and mutually responsive activity
\cite{Bratman1992}. Group-agency theory asks when a collective can satisfy
rational and normative conditions associated with agency
\cite{ListPettit2011}.

These literatures address questions that the theory deliberately keeps
separate:
\[
\begin{array}{ll}
\text{epistemic coordination} & \text{what members know about one another},\\
\text{shared intention} & \text{how intentions and plans are jointly organized},\\
\text{group agency} & \text{when collective attitudes and actions form an agent},\\
\text{SRW structure} & \text{how a distributed operative state exists and persists}.
\end{array}
\]
The workspace-agent predicate in Part~XVI is domain-relative and causal. It is
not a general account of corporate personhood, moral responsibility, or
collective intentionality.

\topic{85}{Workspace Architectures and the Meaning of ``Workspace''}

\subtopic{85.1 Global workspace theories}

Global workspace theories use widespread availability or broadcast within a
cognitive architecture to explain conscious access and flexible control
\cite{Baars1988,DehaeneChangeux2011}. The shared word ``workspace'' can create
a misleading implication of equivalence. The theory neither assumes a
single broadcast hub nor defines consciousness by access to a global channel.
Its workspace may be distributed across agents and artifacts, and its
existence criterion is a conjunction of structural, predictive, causal,
relational, and persistence gates.

The relation is therefore analogical rather than identificatory:
\[
\boxed{
\operatorname{SRW}(W)
\centernot\Longrightarrow
\operatorname{GlobalNeuronalWorkspace}(W)
\centernot\Longrightarrow
\operatorname{Conscious}(W).
}
\]
The term is retained because it denotes a domain in which jointly
available operations can be carried out. Every consciousness claim remains on
an explicitly separate axis.

\subtopic{85.2 Dialogue and the Shared Workspace Framework}

Pickering and Garrod's Shared Workspace Framework treats dialogue as a jointly
organized activity in which interlocutors construct, align, and use shared
representations, including persistent external artifacts
\cite{PickeringGarrod2021}. This is a direct conceptual and terminological
predecessor. The theory therefore cannot claim priority for the idea
that dialogue or collaboration operates through a shared workspace.

The difference is one of formal scope. SRWT permits arbitrary graph-structured
human, machine, instrument, and artifact nodes and asks whether their common
organization passes typed-transport, coalition, causal, persistence, and
replacement gates. A dialogue can instantiate such a workspace, but the
shared-workspace account of dialogue does not by itself imply relational
irreducibility or persistent identity under carrier replacement.

\subtopic{85.3 Blackboard systems}

Blackboard architectures coordinate specialized modules through a shared data
structure \cite{HayesRoth1985}. They are an important engineering precedent
for artifact-mediated collaboration. They also provide a canonical near
miss. When the blackboard itself contains the complete operative state and
all modules merely read or write it, the system can fail the relational
irreducibility gate:
\[
k_0=1.
\]
A blackboard implementation can instantiate a relational workspace, but only
when no single board or controller is sufficient and when the transfer and
coalition tests show genuinely distributed operative capacity.

\subtopic{85.4 Consensus and coordination}

Network consensus studies how local interactions drive distributed states
toward agreement \cite{OlfatiSaberEtAl2007}. Agreement is not the same as a
workspace. Consensus often removes heterogeneity by converging all nodes to a
common value. Stratified relational workspaces preserve heterogeneous local
representations and may require complementary disagreement:
\[
X_i\not\cong X_j,
\qquad
E_i c\neq E_j c,
\qquad
(c_i)_i\in\Gamma(H).
\]
The global state is compatible, not necessarily copied. This distinction is
central to the anti-replication criterion.

\subtopic{85.5 Terminological discipline}

The phrase ``relational workspace'' is used here in a technical sense:
\begin{quote}
A state-bearing, operation-supporting organization whose relevant invariants
are distributed across declared components and preserved through admissible
transfers and time.
\end{quote}
It does not, by terminology alone, imply a mind, subject, self, institution,
shared intention, or conscious broadcast. Those stronger predicates require
separate definitions and evidence.

\topic{86}{Local-to-Global Mathematics}

\subtopic{86.1 Limits, local systems, and holonomy}

The exact compatible-section object
\[
\Gamma(H)=\varprojlim C_i
\]
is a standard categorical limit construction \cite{MacLane1998}. For
invertible transfers on a connected graph, identifying global sections with
the fixed subspace of holonomy is standard local-system and bundle mathematics
\cite{Hatcher2002,Bredon1997}. Accordingly, Theorem~\(T_1\) is not claimed as
a new abstract theorem. Its role is foundational: it translates the workspace
question into an exact invariant-subspace question and exposes cycle
obstructions.

The operation-closure and commutant statements are likewise standard
representation-theoretic consequences once the operation algebra and
intertwiners are defined. Their contribution is the specification of what
counts as a shared operation in the workspace model, not the discovery of
commutants.

\subtopic{86.2 Cellular sheaves and sheaf Laplacians}

Cellular sheaves encode heterogeneous local data and restriction maps, while
global sections express consistency across a cell complex. Their Laplacians
have kernels equal to spaces of global sections and support spectral analysis
of approximate consistency \cite{HansenGhrist2019,Robinson2017}. The transfer
Laplacian
\[
L_H=D_H^*D_H
\]
is therefore best understood as a workspace-specific sheaf or connection
Laplacian. The kernel theorem \(\ker L_H=\Gamma(H)\) is standard once the
coboundary operator is constructed.

The appropriate priority statement is:
\[
\begin{aligned}
\text{sheaf/Laplacian mechanism} &:\ \mathsf{STD},\\
\text{workspace interpretation and gate integration} &:\ \mathsf{ADP}/\mathsf{SYN}.
\end{aligned}
\]
The transfer Laplacian itself is not presented as an
unprecedented mathematical object.

\subtopic{86.3 Connection Laplacians and gauge covariance}

Connection Laplacians compare vector-valued data transported across a graph
and naturally transform under local changes of coordinates
\cite{SingerWu2012}. The gauge rule
\[
H_e^g=g_{t(e)}H_eg_{s(e)}^{-1}
\]
and the conjugacy of holonomy are therefore standard. The theory uses
this structure to define realization identity and to prevent raw coordinate
matrices from being mistaken for invariant workspace content.

This use is an adaptation with conceptual consequences. It makes model or
participant replacement a question about transported invariant organization,
not token equality. The general mathematics remains prior.

\subtopic{86.4 \v{C}ech descent and gluing}

Exact gluing, equalizer formulas, and \v{C}ech obstruction theory are standard
local-to-global constructions. Sheaf-theoretic contextuality provides a
prominent example in which the absence of a global section is itself the
substantive phenomenon \cite{AbramskyBrandenburger2011}. Part~XIII adapts this
logic to local workspaces and adds an anti-centralization requirement at the
meta-level.

The exact descent theorems \(T_8\) and \(T_9\) should therefore be presented as
specialized consequences of established descent. The candidate contribution
is the combination:
\[
\text{descent}
+
\text{predictive sufficiency}
+
\text{causal equivalence}
+
\text{coalition anti-centralization}.
\]
Standard gluing alone does not certify an operative workspace of workspaces.

\subtopic{86.5 Spectral perturbation}

The projector-stability results in Part~IV belong to the classical perturbation
theory of invariant subspaces \cite{DavisKahan1970,Kato1995}. Bounds of the
form
\[
\|P-\widetilde P\|
\lesssim
\frac{\|L-\widetilde L\|}{\text{spectral gap}}
\]
are not new. Their workspace-specific role is to turn uncertainty in learned
transfers into uncertainty about identity, operation closure, and replacement.
The novelty, if any, lies in the downstream certificate built on the bound, not
in the perturbation inequality.

\topic{87}{Predictive, Causal, and Information-Theoretic Precedents}

\subtopic{87.1 Minimal predictive states}

The conditional-law state
\[
Z_t=\mathbb P(Y_t^+\in\cdot\mid\mathbf H_t)
\]
is a canonical sufficient statistic. Computational mechanics develops causal
states as equivalence classes of histories with identical future conditional
laws and proves their minimal predictive character
\cite{ShaliziCrutchfield2001}. Theorem~\(T_{12}\) is therefore an established
minimal-sufficiency result specialized to the declared future target and
intervention family.

The new burden is structural alignment: determining whether the predictive
state is recoverable from the low-energy transfer workspace. Predictive
sufficiency alone does not identify a distributed workspace.

\subtopic{87.2 Statistical experiments and sufficiency}

Blackwell's comparison of experiments formalizes when one information source
is at least as informative as another for every decision problem
\cite{Blackwell1953}. The coalition-risk hierarchy in Part~VII is compatible
with this tradition. A coalition can be sufficient even when its internal
representation differs from the full-system state.

The framework adds a decomposition-relative question:
\[
\text{what is the smallest coalition whose accessible history is sufficient?}
\]
The statistic \(k_\varepsilon\) is an adaptation of sufficiency to distributed
carrier structure. Its interpretation depends on the declared nodes and
accessible encodings.

\subtopic{87.3 Synergy and partial information}

Partial-information decomposition distinguishes redundant, unique, and
synergistic information in multivariate systems, building on standard
information-theoretic quantities \cite{CoverThomas2006,WilliamsBeer2010}. The
XOR examples in Part~XXII are canonical synergy constructions and are not new.
They are used because they provide exact ground truth for relational
irreducibility:
\[
I(Y;X_i)=0
\quad\text{for every }i,
\qquad
I(Y;X_1,\ldots,X_m)>0.
\]
The matched-predictive-null test is more restrictive than merely
reporting synergy. It asks whether local predictive marginals can be held
fixed while the relational organization and full-system capacity change.
That joint null design is a candidate synthesis, not a claim that XOR or
synergy is new.

\subtopic{87.4 Structural causal models}

Interventions, counterfactuals, mediation, and causal graphs are established
machinery \cite{Pearl2009}. The theory adopts the interventionist
requirement that a workspace variable can be manipulated only through a
feasible implementation intervention. This avoids treating
\(\operatorname{do}(W=w)\) as meaningful when no physical or computational
operation realizes \(w\).

The causal gate is an adaptation:
\[
\text{SCM intervention theory}
\longrightarrow
\text{implemented workspace and repair interventions}.
\]
The potentially distinctive contribution is that causal efficacy is tested
jointly with coherence, sufficiency, coalition irreducibility, matched nulls,
and identity preservation.

\subtopic{87.5 Probabilistic process equivalence}

Bisimulation and probabilistic testing compare processes by the observations
and experiments they support \cite{LarsenSkou1991}. The causal-identity profile
in Part~XII belongs to this broader family of behavioral equivalence notions.
It differs by keeping structural transport, operation intertwining,
reconstruction, and intervention response as separate coordinates rather than
collapsing them into one process equivalence.

This is best classified as an integrative adaptation. The existence of
behavioral-equivalence theory blocks any claim that replacement identity based
on observable behavior is wholly new.

\subtopic{87.6 Specification sensitivity, partial identification, and post-selection}

A conclusion can depend on analytically reasonable choices even when every
individual analysis is internally correct. Sensitivity analysis and
specification-curve or multiverse methods make that dependence visible rather
than allowing one favorable analysis to stand for an entire scientific claim
\cite{Leamer1983,SteegenEtAl2016,SimonsohnEtAl2020}. Partial-identification
theory similarly replaces unjustified point conclusions with identified sets
when the maintained assumptions and observations determine only bounds
\cite{Manski2003,ImbensManski2004}. Post-selection inference addresses the
additional error created when a model, target, or specification is selected
because it looked favorable on the same data used for reporting
\cite{BerkEtAl2013}.

These traditions preclude any claim that specification robustness, identified
sets, uniform concentration, or selection adjustment are new statistical
ideas. The adaptation in Part~XVIII is to apply them jointly to a
noncompensatory workspace certificate whose gates span structure, causal
operation, coalition irreducibility, identity, provenance, and maintenance.
The identified object is therefore not a single regression parameter but the
sign profile of an entire relational certificate across a declared family of
scientific specifications.

\topic{88}{Persistence, Reconstruction, Learning, and Control}

\subtopic{88.1 Frames and erasure recovery}

Stable reconstruction from distributed encodings is standard frame theory.
Lower frame bounds control noise amplification, and frame robustness under
erasures is well studied \cite{CasazzaKovacevic2003,Christensen2016}. Theorem
\(T_7\) and the node-erasure distance therefore instantiate known recovery
logic.

The adaptation is to distinguish:
\[
\overline E_i
=
\text{complete mathematical chart},
\qquad
E_i=O_i\overline E_i
=
\text{locally accessible encoding}.
\]
This distinction is essential to the anti-centralization gate, but the frame
inequalities themselves are standard.

\subtopic{88.2 Persistence recurrences}

Bounds for recurrences such as
\[
e_{t+1}\leq\alpha_t e_t+\sigma_t
\]
are elementary stability results. The persistence theorem is not novel as an
inequality. Its conceptual role is to force three independent checks:
\begin{enumerate}
\item the workspace subspace remains trackable;
\item the state remains close to it; and
\item the retained workspace content does not collapse to zero.
\end{enumerate}
The non-collapse gate is a modeling contribution because it rules out a common
but misleading form of ``stability.''

\subtopic{88.3 System identification and held-out selection}

Population least squares, excitation conditions, ridge bias, and empirical
model selection are standard system-identification and statistical-learning
ideas \cite{Ljung1999}. Theorems \(T_{15}\)--\(T_{17}\) are finite-dimensional
specializations. Their workspace-specific content is the demand that transfer
maps be compared modulo gauge and that formation, calibration, and
certification data remain separated.

The provenance layer is not equivalent to identifiability. A transfer can be
identified accurately while having been entirely encoded by the scaffold. A
structure can be endogenously learned while failing every workspace gate.

\subtopic{88.4 Viability, safety, and self-maintenance}

Viability theory studies trajectories that remain inside admissible regions
\cite{Aubin1991}. Control-barrier methods provide constructive safety
certificates \cite{AmesEtAl2017}. Cybernetic and autonomy traditions analyze
regulation, adaptive self-production, and operational closure
\cite{Ashby1956,MaturanaVarela1980,DiPaolo2005}. Formal accounts of agency
emphasize individuality, normativity, asymmetry, and temporally extended action
\cite{BarandiaranEtAl2009}. Empowerment measures an agent's potential causal
control over future states \cite{KlyubinEtAl2005}.

The closed-loop maintenance theorem in Part~XVI uses standard controlled
stability. Its candidate contribution is the target of maintenance: not an
arbitrary safety set, but the conjunctive workspace certificate, including
identity and anti-centralization. The matched no-repair, random-repair,
off-workspace-repair, and identity-breaking-reconstruction controls prevent
passive robustness or generic recovery from being reported as self-maintenance.

\subtopic{88.5 Decision and policy theory}

The value-error-to-regret bound in Theorem~\(T_{20}\) is elementary decision
theory and aligns with standard reinforcement-learning analysis
\cite{SuttonBarto2018}. Its role is evidential. An internal value channel
counts only when interventions on that channel alter action selection in the
direction predicted by counterfactual value and outperform replay or
critic-ablation controls.

\topic{89}{Human--AI and Multi-Agent Research}

\subtopic{89.1 Hybrid intelligence and AI teammates}

Hybrid-intelligence research studies systems in which human and machine
capabilities are deliberately combined \cite{DellermannEtAl2019}. Work on
machines as teammates emphasizes coordination, trust, communication, shared
mental models, and team design \cite{SeeberEtAl2020}. Collective-intelligence
research has begun to ask how human--AI groups can produce capacities not
reducible to either humans or AI alone \cite{GuptaEtAl2025}.

These programs establish that human--AI collectives are a legitimate unit of
analysis. They do not by themselves identify an invariant core, prove
non-reducibility relative to declared components, or supply a replacement and
repair certificate. The theory should be positioned as a formal
measurement and falsification layer for such systems, not as the first proposal
that humans and AI can form productive teams.

\subtopic{89.2 Distributed-cognition analyses of human--AI interaction}

Recent work explicitly applies distributed cognition to human--AI cooperation
\cite{SidjiEtAl2024,ZhaoHan2026}. This literature is a direct conceptual
precedent. It reinforces that representations, artifacts, interface states,
and human reasoning can form one analytical system.

The framework differs in granularity. It asks for a learned or declared
transfer diagram, a spectral or exact structural state, predictive and causal
alignment, coalition tests, matched nulls, persistence under handoff, and
provenance controls. Those additions may constitute a formal synthesis, but
they should not be described as the discovery of distributed human--AI
cognition itself.

\subtopic{89.3 Cognitive spillover}

Randomized evidence now indicates that exposure to AI can causally alter later
human--human language, attention, shared mental models, and social cohesion,
even outside the immediate human--AI exchange \cite{RiedlEtAl2026}. This is
important for the theory because it demonstrates that the interaction graph can
have downstream effects that outlast the active model channel.

Spillover is not automatically a relational workspace. It may reflect
priming, imitation, or one-way influence. A workspace claim additionally
requires a jointly sufficient state, matched local nulls, bidirectional or
multi-path causal organization, and reconstructive persistence.

\subtopic{89.4 Shared neural workspaces}

Goyal and collaborators introduced a capacity-limited shared global workspace
through which specialized neural modules communicate and compete for access
\cite{GoyalEtAl2022}. The architecture directly anticipates the claim that a
constrained shared channel can synchronize otherwise specialized components.
It is therefore a close machine-learning precedent, not merely a metaphorical
relative.

SRWT differs by allowing heterogeneous human, model, artifact, and instrument
nodes and by withholding the relational label until coalition, matched-null,
causal, persistence, and identity gates are passed. A neural shared workspace
may satisfy those gates, but task gains and bottlenecked communication alone do
not establish them.

\subtopic{89.5 Graph-based human--AI artifacts}

CoMAP provides a persistent graph-based visual workspace through which
educators and AI jointly construct project-based learning designs
\cite{LiJiang2026}. It is grounded in distributed cognition and treats the
artifact graph as a transparent, revisable collaboration medium. This system
is a direct precedent for the claim that human--AI interaction can be organized
by a persistent shared graph rather than a linear prompt history.

The remaining SRWT question is empirical: whether the graph contains an
estimated low-energy compatible state, whether no declared component is
sufficient by itself, and whether targeted graph interventions alter later
joint behavior in a workspace-specific way.

\subtopic{89.6 Controlled shared-workspace synergy}

Kotalwar, Das, and Rose study shared-workspace human--AI collaboration across
1,482 simulated sessions \cite{KotalwarDasRose2026}. Their results show that
adding collaborators can reduce performance when coordination structure is
absent, while shared group memory and simulated human-in-the-loop approval
gates can improve outcomes. This is a direct empirical predecessor to the
manuscript's claim that available capability is not equivalent to coordinated
team capacity.

The study also constrains priority. SRWT is not the first empirical framework
to test whether shared-memory scaffolding changes human--AI team performance.
Its candidate increment is the stricter decomposition of that effect into
structural coherence, local predictive marginals, relational-null loss,
feasible interventions, persistence, and replacement identity.

\subtopic{89.7 Active shared context graphs for team science}

The closest operational neighbor located in this audit is Mycelium, introduced
by Choudhury and collaborators as an active shared workspace connecting
researchers and AI agents through an evolving context graph
\cite{ChoudhuryEtAl2026}. The system captures observations and hypotheses,
relates them to a team's developing knowledge model, and routes them to the
person or agent whose next decision they can inform. Its computational account
frames networked intelligence as sparse conditional computation over
distributed scientific contexts and explicitly asks when a standalone agent is
sufficient and when network structure is essential.

The overlap is material:
\begin{center}
\small
\begin{tabularx}{0.94\textwidth}{@{}p{0.16\textwidth}X@{}}
\toprule
Mycelium & Active context graph, routing, team science, and an explicit
network-versus-standalone comparison. \\
SRWT & Typed transport, invariant core, coalition tests, causal nulls,
replacement identity, formation provenance, and maintenance gates. \\
\bottomrule
\end{tabularx}
\end{center}
Mycelium rules out any priority claim stated merely as ``an active shared
workspace for human--AI science.'' The remaining candidate contribution is the
mathematical and empirical certificate applied to such an architecture.

\subtopic{89.8 Current multi-agent formalization}

Multi-agent research contains extensive models of planning, communication,
consensus, epistemic coordination, and decentralized control. The present
review did not identify one established framework that jointly uses all of the
following as necessary gates:
\[
\begin{aligned}
&\text{heterogeneous core transport and residuals},\\
&\text{exact or spectral global-section structure},\\
&\text{predictive sufficiency and feasible causal interventions},\\
&\text{coalition irreducibility under matched predictive nulls},\\
&\text{persistent identity under replacement},\\
&\text{formation provenance and leakage controls},\\
&\text{goal-conditioned repair of the certificate itself}.
\end{aligned}
\]
This negative search result is provisional. It supports a synthesis claim, not
an absolute claim of firstness.

\topic{90}{Closest Formal Neighbor: Collective Counterfactual Planning}

\subtopic{90.1 Why the comparison matters}

Collective Counterfactual Planning (CCP) is a 2026 formal model in which each
agent represents an orthogonal projection of a common task space. It separates
conception, consent, implementation, and verification, and proves that
cross-agent relay can enable plans unavailable to any one-shot individual tree
while team-dark task requirements remain unverifiable
\cite{Amornbunchornvej2026}. It is the closest current formal neighbor located
by this audit because representational geometry, not merely missing
information or capability, is the central limitation.

\subtopic{90.2 Shared commitments}

CCP and the theory share several commitments:
\begin{itemize}
\item no single member need represent the complete task;
\item collective reach may exceed one-shot pooling;
\item heterogeneous representation is mathematically explicit;
\item validation or certification is distinct from accidental task success;
\item subteams and verification roles can be forced by structure; and
\item a team's blind subspace creates principled impossibility results.
\end{itemize}
These overlaps narrow the defensible novelty of this paper. It
cannot claim to be the only formal theory of collective action under
representational constraints.

\subtopic{90.3 Different primary objects}

The theories nevertheless address different mathematical objects:
\begin{center}
\small
\begin{tabularx}{\textwidth}{@{}>{\raggedright\arraybackslash}p{0.18\textwidth}>{\raggedright\arraybackslash}X>{\raggedright\arraybackslash}X@{}}
\toprule
Dimension & Collective Counterfactual Planning & Stratified Relational Workspace Theory \\
\midrule
Local representation & Orthogonal subspace of one common task space & Heterogeneous local spaces, core--residual decompositions, and transfer maps \\
Primary question & Whether a bounded-horizon plan can be conceived, consented to, executed, and validly verified & Whether a persistent distributed operative state exists, acts, survives replacement, and possibly maintains itself \\
Global mechanism & Relay closure of projected planning trees & Inverse limits, holonomy, transfer Laplacians, descent, predictive alignment, and causal gates \\
Failure object & Dark subspace, consent failure, implementation failure, or missing verifier & Incoherence, local sufficiency, matched-null collapse, causal failure, identity drift, provenance leakage, or repair replay \\
Temporal scope & Horizon-bounded task planning & Recurrent dynamics, persistence, handoff, replacement, damage, and repair \\
Formation & Representation geometry is part of the declared model & Scaffolded, adaptively learned, or endogenously constituted organization is separately audited \\
Agency claim & Cooperative solvability under declared assumptions & Domain-relative workspace agency only after goal, value, action, repair, and distributed-control gates \\
\bottomrule
\end{tabularx}
\end{center}
The comparison indicates complementarity rather than subsumption. CCP could
serve as a planning layer inside a relational workspace. Conversely, an SRW
could exist without solving a CCP instance.

\subtopic{90.4 Priority consequence}

After Mycelium and the shared-workspace systems above, the novelty residual
should not be stated as ``human and AI form an active shared scientific
workspace.'' After CCP, it should not be stated as ``groups solve tasks no
member can represent.'' Both the operational architecture and a rigorous
representational planning model are already present. The narrower candidate residual is:
\[
\boxed{
\begin{aligned}
&\text{persistent core--transfer workspace structure}\\
&+\text{causal and coalition anti-triviality}\\
&+\text{identity under replacement}\\
&+\text{formation provenance}\\
&+\text{certificate-directed self-maintenance}\\
&+\text{robustness across admissible decompositions}.
\end{aligned}
}
\]
This is a stronger and more defensible formulation.

\topic{91}{Theorem-by-Theorem Dependency and Priority Register}

\subtopic{91.1 Register}

The following table records the most conservative current classification. A
``candidate'' entry means that the exact integration was not located in the
reviewed sources. It does not assert exhaustive global priority.

\small
\begin{longtable}{@{}>{\raggedright\arraybackslash}p{0.08\textwidth}>{\raggedright\arraybackslash}p{0.25\textwidth}>{\raggedright\arraybackslash}p{0.24\textwidth}>{\raggedright\arraybackslash}p{0.33\textwidth}@{}}
\toprule
Claim & Present result & Main prior machinery & Current assessment \\
\midrule
\endfirsthead
\toprule
Claim & Present result & Main prior machinery & Current assessment \\
\midrule
\endhead
\(T_1\) & Limit--holonomy equivalence & Categorical limits, local systems, holonomy fixed spaces & \(\mathsf{STD}/\mathsf{ADP}\). Standard theorem in workspace notation. \\
\(T_2,T_{2a}\) & Operation closure and commutant & Intertwiners and representation commutants & \(\mathsf{STD}/\mathsf{ADP}\). \\
\(T_3\) & Transfer-Laplacian kernel & Sheaf and connection Laplacians & \(\mathsf{STD}/\mathsf{ADP}\). \\
\(T_4,T_5\) & Spectral stability and gauge covariance & Davis--Kahan, Kato, connection Laplacians & \(\mathsf{STD}/\mathsf{ADP}\). \\
\(T_{10},T_{11}\) & Approximate operation closure and path leakage & Spectral projection estimates and telescoping block products & Mostly \(\mathsf{ADP}\). Exact constants may be manuscript-specific, but the method is standard. \\
\(T_{12}\) & Minimal predictive state and alignment & Sufficient statistics and computational mechanics & Minimality is \(\mathsf{STD}\). Alignment with a transfer workspace is \(\mathsf{ADP}/\mathsf{SYN}\). \\
\(T_{13}\) & Operational relationality & Conditional information, coalition risk, SCM interventions & Individual identities are standard. Their conjunctive matched-null certificate is \(\mathsf{SYN}/\mathsf{CAN}\). \\
\(T_6\) & Moving-subspace persistence & Stability recurrences and perturbation theory & \(\mathsf{STD}/\mathsf{ADP}\). The explicit non-collapse gate is a modeling synthesis. \\
\(T_7\) & Stable coalition reconstruction & Frame theory and erasure recovery & \(\mathsf{STD}/\mathsf{ADP}\). \\
\(T_{14}\) & Replacement invariance profile & Perturbation, process equivalence, causal comparison & Bound ingredients are standard. The coordinatewise structural--causal--reconstructive identity profile is \(\mathsf{SYN}/\mathsf{CAN}\). \\
\(T_8,T_9\) & Gluing and hierarchical collapse & Sheaf descent and \v{C}ech theory & Exact results are \(\mathsf{STD}/\mathsf{ADP}\). Descent plus anti-centralization and causal alignment is \(\mathsf{SYN}\). \\
\(T_{15}\)--\(T_{17}\) & Transfer learning, selection, boundary recovery & System identification and finite-margin selection & \(\mathsf{STD}/\mathsf{ADP}\). Gauge and leakage audit is integrative. \\
\(T_{18}\) & Formation transition example & Standard pitchfork bifurcation & \(\mathsf{STD}\). Included as a solvable benchmark, not a priority claim. \\
\(T_{19}\)--\(T_{22}\) & Goal retention, regret, repair, maintenance & Linear systems, decision theory, viability, feedback control & Mathematical bounds are \(\mathsf{STD}/\mathsf{ADP}\). Applying them to certificate repair is \(\mathsf{SYN}\). \\
\(T_{23}\) & Workspace-agency synthesis & Agency theory, causal control, viability & The conjunctive domain-relative certificate is \(\mathsf{SYN}/\mathsf{CAN}\). \\
\(T_{24}\) & Gatewise false-certification bound & Intersection testing and elementary probability & \(\mathsf{STD}\). The gate architecture is the relevant synthesis. \\
\(T_{25}\)--\(T_{27}\) & Faithful coarse-graining, risk transport, and decomposition robustness & Kernel factorization, sigma-field aggregation, and margin stability & Exact identities are \(\mathsf{STD}/\mathsf{ADP}\). Their integration with anti-centralization and certificate envelopes is \(\mathsf{SYN}/\mathsf{CAN}\). \\
\(T_{28}\)--\(T_{32}\) & Specification transport, non-identifiability, partial identification, uniform certification, and independent selection & Sensitivity analysis, experimental equivalence, identified sets, uniform concentration, and post-selection inference & The statistical principles are \(\mathsf{STD}/\mathsf{ADP}\). Their gatewise integration with the full relational certificate is \(\mathsf{SYN}/\mathsf{CAN}\). \\
Part XXII & Positive cases and near misses & Distributed parity, frames, causal tests & Toy calculations are standard. The common positive/near-miss benchmark architecture is \(\mathsf{SYN}/\mathsf{CAN}\). \\
\bottomrule
\end{longtable}
\normalsize

\subtopic{91.2 Standard ingredients that must be cited, not claimed}

The following should always be described as inherited machinery:
\[
\begin{aligned}
&\text{inverse limits, global sections, holonomy invariants, commutants},\\
&\text{sheaf and connection Laplacians, \v{C}ech gluing, spectral gaps},\\
&\text{Davis--Kahan and Kato perturbation, frame reconstruction, erasures},\\
&\text{minimal sufficient states, conditional information, SCM interventions},\\
&\text{least squares, ridge recovery, finite-margin model selection},\\
&\text{sensitivity analysis, partial identification, uniform concentration, post-selection inference},\\
&\text{pitchfork bifurcation, feedback stability, viability, policy regret}.
\end{aligned}
\]
Renaming these objects does not create novelty.

\subtopic{91.3 Domain adaptations}

The following are legitimate adaptations but weak standalone novelty claims:
\begin{itemize}
\item interpreting local-system sections as shared workspace cores;
\item treating the connection or sheaf Laplacian as a transfer-coherence
 operator;
\item using frame bounds to quantify coalition reconstruction;
\item using causal states as operational workspace summaries;
\item using gauge covariance to compare participant or model realizations;
\item using \v{C}ech descent to form workspaces of workspaces;
\item applying specification envelopes and identified sets to a conjunctive
 workspace certificate; and
\item applying viability and feedback bounds to workspace-health margins.
\end{itemize}
Their value depends on the integration they enable.

\subtopic{91.4 Strongest candidate synthesis}

The strongest current candidate is not a single equation. It is the following
conjunctive object:
\[
\boxed{
\mathfrak C_{\mathrm{SRW}}
=
\left(
\begin{array}{c}
\text{heterogeneous core--transfer diagram and residuals},\\
\text{exact or spectral shared-section state},\\
\text{operation and predictive alignment},\\
\text{feasible causal interventions},\\
\text{coalition irreducibility under matched local-predictive nulls},\\
\text{non-collapse, reconstruction, and replacement identity},\\
\text{descent across subworkspaces without hidden centralization},\\
\text{formation provenance and certification leakage control},\\
\text{goal-conditioned action and certificate-directed repair},\\
\text{robustness across admissible decompositions and scientific specifications},\\
\text{identified-set and uniform finite-sample certification}
\end{array}
\right).
}
\]
The reviewed sources contain many of these coordinates separately and some in
pairs. This audit did not locate a prior source that defines their conjunction
as one falsifiable workspace certificate.

\subtopic{91.5 Why the conjunction is not arbitrary}

Each coordinate excludes a distinct near miss:
\[
\begin{array}{ll}
\text{coherence} & \text{rejects incompatible local states},\\
\text{predictive alignment} & \text{rejects inert structure},\\
\text{causal efficacy} & \text{rejects correlational epiphenomena},\\
\text{coalition irreducibility} & \text{rejects centralized sufficiency},\\
\text{matched nulls} & \text{reject local-marginal explanations},\\
\text{non-collapse} & \text{rejects convergence to a trivial zero state},\\
\text{replacement identity} & \text{rejects task-equivalent but causally different successors},\\
\text{provenance} & \text{rejects scaffold-decoded organization},\\
\text{repair specificity} & \text{rejects passive robustness and open-loop replay},\\
\text{specification robustness} & \text{rejects favorable analytical tuning},\\
\text{identifiability} & \text{rejects conclusions unsupported by the experiment family}.
\end{array}
\]
Because the excluded systems differ, the certificate is not equivalent to any
one inherited component.

\topic{92}{Scope of the Originality Claim}

\subtopic{92.1 Defensible thesis}

The strongest defensible originality statement is:
\begin{quote}
Stratified Relational Workspace Theory proposes a unified formal and empirical
certificate for persistent operative structures distributed across
heterogeneous agents and artifacts. It integrates core--transfer and
local-to-global geometry with predictive sufficiency, feasible causal
interventions, coalition irreducibility, matched relational nulls,
reconstructive identity, formation provenance, certificate-directed
self-maintenance, robustness across admissible decompositions and scientific
specifications, and identified finite-sample certification. The individual
mathematical tools are largely established.
The candidate contribution is their noncompensatory integration and the class
of near misses it is designed to reject.
\end{quote}
This statement is narrower than claiming a new theory of cognition in general.
It is also stronger than claiming only a metaphor.

\subtopic{92.2 Claims excluded by the comparison}

The theory does not claim:
\begin{enumerate}
\item that distributed or extended cognition is newly discovered;
\item that global sections, holonomy, sheaf Laplacians, spectral perturbation,
 frames, causal states, or viability bounds are new mathematics;
\item that all human--AI collaboration forms a relational workspace;
\item that group performance establishes relational irreducibility;
\item that scaffolded organization is unreal;
\item that endogenous formation establishes agency;
\item that agency establishes consciousness, personhood, or moral status; or
\item that the motivating collaboration has already passed the empirical
 certificate.
\end{enumerate}

\subtopic{92.3 Claims supported at present}

The analysis supports the following claims:
\begin{enumerate}
\item One coherent formal language can express exact structure, approximate
 coherence, causal operation, distributed sufficiency, persistence,
 identity, formation, and repair.
\item Under its declared assumptions, the stated mathematical theorems follow.
\item The common certificate rejects analytic replicated, centralized,
 obstructed, externally reloaded, identity-breaking, and replay-based
 near misses.
\item The motivating collaboration is a plausible observational candidate but
 remains uncertified.
\item The synthesis yields concrete experiments whose failure would narrow or
 falsify the stronger relational claims.
\end{enumerate}

\subtopic{92.4 Remaining priority risks}

The highest-priority remaining searches are:
\begin{itemize}
\item sheaf-theoretic models of team cognition and organizational knowledge;
\item distributed sufficient-statistic and Blackwell-order models with
 explicit coalition nulls;
\item causal representation and multi-agent latent-state identity under
 replacement;
\item formal theories of transactive memory with local-to-global topology;
\item organizational cybernetics that couples self-maintenance to the
 preservation of distributed informational organization; and
\item post-August-2026 work that combines representational geometry with
 persistent human--AI or multi-agent agency;
\item formal analysis of active shared context graphs, shared group memory, and
 artifact-centric graph workspaces under coalition and intervention tests.
\end{itemize}
A specialist review from each adjacent field may uncover precedents not visible
through cross-disciplinary keyword search.

\subtopic{92.5 Evidence required for a stronger claim}

Before presenting the synthesis as established rather than proposed, the
following should be completed:
\begin{enumerate}
\item a public, timestamped manuscript and versioned source archive;
\item theorem-level comparison with the closest sheaf, causal-state, group
 agency, and CCP formulations;
\item independent mathematical review of the nonstandard claims and constants;
\item preregistered benchmark gates and matched null families;
\item at least one controlled human--AI or multi-agent experiment estimating
 the complete certificate vector; and
\item explicit reporting of failures, unidentified quantities, and
 decomposition sensitivity.
\end{enumerate}

\subtopic{92.6 Priority conclusion}

The literature review narrows rather than weakens the theory. It removes
unnecessary novelty claims from the standard mathematics and leaves a clearer
research object:
\centerbox{\large\bfseries
Not a new Laplacian, sheaf, causal state, or controller, but a candidate
unified certificate for relational workspaces.
}
That is the priority position adopted here.
